\subsection{Prospective class-conditioned assay calibration}
A separate prospective pilot used five original seeds and 15 exact-retraining seeds
per condition, without evaluating approximate methods. Full deletion of Fashion-MNIST
class 0 and class 5 produced JS-divergence separability ratios of
\PhaseThreeClassZeroJSRatio{} and \PhaseThreeClassFiveJSRatio{}, exceeding their
simulated minimum detectable ratios of \PhaseThreeClassZeroMinimum{} and
\PhaseThreeClassFiveMinimum{}. Half-class deletion and every tested disagreement,
CKA-distance, and MMD endpoint remained assay-insensitive. A model-level precision
simulation selected ten exact seeds for JS under the frozen 20\% interval rule in both
full-class conditions.

\subsection{Prospective equivalence-margin calibration}
A second frozen pilot tested whether the condition-specific exact/exact 95th percentile
could detect a controlled violation in which 1\% of the forbidden class remained in
training. The candidate margins were \PhaseThreeClassZeroMargin{} for class 0 and
\PhaseThreeClassFiveMargin{} for class 5. Detection probabilities were only
\PhaseThreeClassZeroDetection{} and \PhaseThreeClassFiveDetection{}, respectively,
below the prospectively required 80\% in both cases. The candidate margins are therefore
inadmissible. Larger contamination levels were designated diagnostics and cannot rescue
the failed primary rule. No approximate-method equivalence claim is made.
